\begin{table}[H]
\caption{EE Mobility and Workers’ Employment Expectations}\label{tab:C_sce}
\fontsize{10.5}{12}\selectfont\noindent
\begin{tabular}{@{}L{108pt}C{84pt}C{76pt}C{64pt}C{80pt}C{56pt}@{}}
\toprule
\multicolumn{6}{@{}l@{}}{\textit{Panel A: Perceived job-finding probability}}\\[3pt]
$h$ (months) & $\hat\beta_h$ & $t_{NW}$ & $p$ & $R^2$ (\%) & $N$\\
\midrule
Current & 0.00 & 0.00 & 0.000 & 0.00 & 0\\[2pt]
1 & 0.00 & 0.00 & 0.000 & 0.00 & 0\\[2pt]
3 & 0.00 & 0.00 & 0.000 & 0.00 & 0\\[2pt]
6 & 0.00 & 0.00 & 0.000 & 0.00 & 0\\[2pt]
12 & 0.00 & 0.00 & 0.000 & 0.00 & 0\\[2pt]
\midrule
\multicolumn{6}{@{}l@{}}{\textit{Panel B: Perceived job-loss probability}}\\[3pt]
$h$ (months) & $\hat\beta_h$ & $t_{NW}$ & $p$ & $R^2$ (\%) & $N$\\
\midrule
Current & 0.00 & 0.00 & 0.000 & 0.00 & 0\\[2pt]
1 & 0.00 & 0.00 & 0.000 & 0.00 & 0\\[2pt]
3 & 0.00 & 0.00 & 0.000 & 0.00 & 0\\[2pt]
6 & 0.00 & 0.00 & 0.000 & 0.00 & 0\\[2pt]
12 & 0.00 & 0.00 & 0.000 & 0.00 & 0\\[2pt]
\bottomrule\end{tabular}
\par\vspace{7pt}
\begin{minipage}[t]{\linewidth}\exhibitnotestyle
\setlength{\abovedisplayskip}{3pt}\setlength{\belowdisplayskip}{3pt}
\setlength{\abovedisplayshortskip}{3pt}\setlength{\belowdisplayshortskip}{3pt}
\textit{Notes:} This table reports ordinary least squares estimates of the relation between EE mobility and employment expectations among currently employed respondents to the New York Fed Survey of Consumer Expectations. The contemporaneous and conditional models are
\[y_t=\alpha+\beta z_t^{EE}+\varepsilon_t,
\qquad
y_{t+h}=\alpha_h+\beta_h z_t^{EE}+\gamma_h y_t+\varepsilon_{t,h},\]
respectively, where $z_t^{EE}$ is standardized EE mobility in economic month $t$. Job finding is the mean perceived probability of finding and accepting a new job within three months if the current job were lost today; job loss is the mean perceived probability of losing the current job within twelve months. Survey probabilities remain in percentage points. The $h=0$ rows report contemporaneous relationships, while $h=1,3,6,12$ indexes the future survey month and the regressions control for the current assessment. These regression horizons do not change the three-month job-finding or twelve-month job-loss event definitions. Observations span June 2013--December 2025. Coefficients measure changes in the reported probability, in percentage points, associated with a one-standard-deviation increase in EE. The $t_{NW}$ column reports Newey--West statistics with 12 monthly lags; $p$ and stars use two-sided Student-$t$ probabilities. $R^2$ is reported in percent, and $N$ is the number of observations. $^{*}$, $^{**}$, and $^{***}$ denote significance at the 10\%, 5\%, and 1\% levels. Appendix \ref{app:employment-security} provides the survey and regression details.\par\end{minipage}
\end{table}
